\begin{longtable}{lllllllrrllllllll}
\toprule
family & system & split & qa\_id & q\_type & dimension & doc\_id & recall\_gain & reworded\_also\_correct & question & reference & base\_answer & base\_verdict & base\_recall & ft\_answer & ft\_verdict & ft\_recall \\
\midrule
\endfirsthead
\toprule
family & system & split & qa\_id & q\_type & dimension & doc\_id & recall\_gain & reworded\_also\_correct & question & reference & base\_answer & base\_verdict & base\_recall & ft\_answer & ft\_verdict & ft\_recall \\
\midrule
\endhead
\midrule
\multicolumn{17}{r}{Continued on next page} \\
\midrule
\endfoot
\bottomrule
\endlastfoot
qwen3-8b & qwen3-8b\_\_ep3 & test\_trained\_exact & ai-roles:0017:q4 & definition & monitoring\_improvement & ai-roles & 1.000000 & True & In the evaluation of an AI early warning system, what are sensitivity (recall) and precision? & Sensitivity (recall) is the fraction of students who actually failed who were correctly predicted to fail; in the hypothetical example it is 60\%. Precision is the fraction of students predicted to fail who actually failed; in the example it is 37.5\%. & In the evaluation of an AI early warning system, **sensitivity (also known as recall)** and **precision** are key performance metrics that measure different aspects of the system's effectiveness: - **Sensitivity (Recall):** Measures the proportion of actual positive cases (e.g., students at risk) that are correctly identified by the system. It reflects the system's ability to detect true positives. **Formula:** \$\$ \textbackslash text\{Sensitivity\} = \textbackslash frac\{\textbackslash text\{True Positives\}\}\{\textbackslash text\{True Positives\} + \textbackslash text\{False Negatives\}\} \$\$ - … [121 words] & partial & 0/4 & Sensitivity (recall) is the fraction of students who actually failed who were correctly predicted to fail; in the hypothetical example it is 60\%. Precision is the fraction of students predicted to fail who actually failed; in the example it is 37.5\%. [41 words] & correct & 4/4 \\
qwen3-8b & qwen3-8b\_\_ep3 & test\_trained\_exact & language-teachers-ai-literacy-a-psychometric:0019:q0 & factual & faculty\_readiness & language-teachers-ai-literacy-a-psychometric & 1.000000 & True & Which six dimensions of AI literacy does the TAILS, a teacher AI literacy scale, measure in preservice English language teachers? & The TAILS measures AI literacy across six dimensions: Knowledge, Evaluation, Collaboration, Contextualization, Autonomy, and Ethics. & The TAILS (Teacher AI Literacy Scale) measures six dimensions of AI literacy in preservice English language teachers: **Technology**, **Awareness**, **Impact**, **Learning**, **Skills**, and **Ethics**. [24 words] & incorrect & 0/1 & The TAILS measures AI literacy across six dimensions: Knowledge, Evaluation, Collaboration, Contextualization, Autonomy, and Ethics. [15 words] & correct & 1/1 \\
qwen3-8b & qwen3-8b\_\_ep3 & test\_trained\_exact & artificial-intelligence-and-accounting-education:0137:q0 & application & teaching\_learning & artificial-intelligence-and-accounting-education & 1.000000 & True & An accounting department wants to introduce generative AI tools for students and is deciding which assistance formats to offer and what concerns to expect. What did a survey of accounting students find about preferred formats and top concerns? & Step-by-step guidance was the most preferred format (45 students), followed by customized quizzes tailored to weak spots (30 students). Few students chose fancier options like visual summaries or real-time fixes, showing a preference for practical, hands-on help. Data privacy topped the list of concerns, with 60 students (33\% of responses) worried about who sees their information. & A survey of accounting students found that preferred formats for generative AI tools include **interactive problem-solving assistance**, **real-time feedback on assignments**, and **practice problem generation**. Students valued tools that could help with **conceptual understanding**, **problem-solving**, and **grammar/structure improvements** in written work. Top concerns included **plagiarism risks**, **reliance on AI leading to reduced learning**, **accuracy and reliability of AI-generated content**, and **ethical concerns around AI use in academic settings**. Students also … [84 words] & partial & 0/4 & Step-by-step guidance was the most preferred format (45 students), followed by customized quizzes tailored to weak spots (30 students). Few students chose fancier options like visual summaries or real-time fixes, showing a preference for practical, hands-on help. Data privacy topped the list of concerns, with 60 students (33\% of responses) worried about who sees their information. [56 words] & correct & 4/4 \\
gemma-4-e4b-it & gemma-4-e4b-it\_\_ep3 & test\_trained\_exact & ai-and-security-in-china-s-higher-education:0070:q1 & explanation & student\_ai\_literacy & ai-and-security-in-china-s-higher-education & 0.800000 & True & Why is there a need for an AI literacy framework designed specifically for higher education? & Most existing frameworks target general or K-12 audiences and under-specify the epistemic demands, disciplinary nuance, and socio-technical complexity of AI use in universities. This gap points to the need for a higher education-oriented framework. & There is a critical need for an AI literacy framework in higher education for several interconnected reasons: 1. **Navigating AI as a Tool and Subject:** Students and faculty are increasingly interacting with AI (e.g., ChatGPT, Copilot) both as a productivity tool and as a subject of study. A framework ensures users understand *how* these tools work, their capabilities, and their limitations, moving beyond mere usage to critical engagement. 2. **Ethical … [213 words] & partial & 1/5 & Most existing frameworks target general or K-12 audiences and under-specify the epistemic demands, disciplinary nuance, and socio-technical complexity of AI use in universities. This gap points to the need for a higher education-oriented framework. [34 words] & correct & 5/5 \\
gemma-4-e4b-it & gemma-4-e4b-it\_\_ep3 & test\_trained\_exact & ai-and-digital-transformation:0039:q4 & application & leadership\_strategy & ai-and-digital-transformation & 1.000000 & True & A small university IT unit wants to adopt the FRE-AIDT ethical AI framework across several campuses in different countries. What challenges of FRE-AIDT might it face? & FRE-AIDT faces challenges including ensuring small organizations can implement complex tools, defining fairness and inclusivity across cultural contexts, and scaling continuous monitoring globally. It also faces the challenge of balancing innovation with compliance costs. & Adopting the FRE-AIDT framework across multiple campuses in different countries presents several significant challenges, primarily related to **contextualization, governance, and implementation consistency.** Here are the key challenges: 1. **Cultural and Legal Divergence (Contextualization):** * **Data Privacy Laws:** Different countries have vastly different data protection regulations (e.g., GDPR in Europe vs. various state laws in the US). The "Fairness" and "Accountability" components of FRE-AIDT must be interpreted and implemented to comply … [291 words] & partial & 0/4 & FRE-AIDT faces challenges including ensuring small organizations can implement complex tools, defining fairness and inclusivity across cultural contexts, and scaling continuous monitoring globally. It also faces the challenge of balancing innovation with compliance costs. [34 words] & correct & 4/4 \\
gemma-4-e4b-it & gemma-4-e4b-it\_\_ep3 & test\_trained\_exact & factors-affecting-performance-expectancy-and:0000:q4 & factual & general & factors-affecting-performance-expectancy-and & 1.000000 & True & Which factor showed a highly significant effect on performance expectancy from AI chatbots in a SmartPLS4 study of ChatGPT-like tools? & Source trustworthiness showed a highly significant effect on performance expectancy from AI chatbots. & To answer this accurately, I would need to **see the specific SmartPLS4 study** you are referring to. Different studies investigating ChatGPT-like tools will yield different results based on their specific model, sample, and context. **In general, when looking at performance expectancy in AI chatbot studies, the most commonly highly significant factors often relate to:** 1. **Perceived Usefulness/Task Relevance:** How well the user believes the chatbot helps them achieve their academic … [110 words] & incorrect & 0/1 & Source trustworthiness showed a highly significant effect on performance expectancy from AI chatbots. [13 words] & correct & 1/1 \\
llama-3.1-8b-instruct & llama-3.1-8b-instruct\_\_ep2 & test\_trained\_exact & ai-and-security-in-china-s-higher-education:0087:q3 & explanation & privacy\_security & ai-and-security-in-china-s-higher-education & 1.000000 & True & What did Rana et al. (2024) investigate about trust in AI, and what insight do their findings offer? & Rana et al. (2024) investigated how trust in AI technologies, shaped by perceptions of transparency and reliability, along with data privacy concerns, influences AI adoption in educational environments. Their findings offer critical insights for encouraging AI integration, especially in resource-limited academic settings. & I cannot verify the findings of Rana et al. (2024) on trust in AI. [14 words] & incorrect & 0/5 & Rana et al. (2024) investigated how trust in AI technologies, shaped by perceptions of transparency and reliability, along with data privacy concerns, influences AI adoption in educational environments. Their findings offer critical insights for encouraging AI integration, especially in resource-limited academic settings. [42 words] & correct & 5/5 \\
llama-3.1-8b-instruct & llama-3.1-8b-instruct\_\_ep2 & test\_trained\_exact & ai-frameworks:0019:q4 & factual & governance & ai-frameworks & 1.000000 & True & What is one instance of an emergent capability revealed when evaluating a newly trained AI model that could lead to harm if left unchecked? & A model capable of autonomously solving CAPTCHA challenges can be used by bad actors to overcome spam protections. Evaluating a newly trained model can reveal such emergent capabilities, which could lead to harm if left unchecked. & One instance of an emergent capability that could lead to harm if left unchecked is the ability of a language model to generate persuasive and manipulative content, such as propaganda or disinformation. For example, a newly trained AI model may be able to generate convincing arguments or persuasive text that can be used to manipulate public opinion or influence people's decisions. If left unchecked, this capability could be used to … [128 words] & incorrect & 0/1 & A model capable of autonomously solving CAPTCHA challenges can be used by bad actors to overcome spam protections. Evaluating a newly trained model can reveal such emergent capabilities, which could lead to harm if left unchecked. [36 words] & correct & 1/1 \\
llama-3.1-8b-instruct & llama-3.1-8b-instruct\_\_ep2 & test\_trained\_exact & responsible-ai-in-the-enterprise-practical-ai:0088:q2 & factual & procurement\_technology & responsible-ai-in-the-enterprise-practical-ai & 0.750000 & True & Which features does IBM's AI Explainability toolkit offer to enhance model explainability? & IBM's AI Explainability toolkit is one of the most mature offerings on the market. It offers global and local explainability, fairness indicators, and counterfactual explanations. & IBM's AI Explainability 360 toolkit offers several features to enhance model explainability, including: 1. **SHAP (SHapley Additive exPlanations)**: A technique to assign a value to each feature for a specific prediction, explaining how much it contributed to the outcome. 2. **LIME (Local Interpretable Model-agnostic Explanations)**: A technique to generate an interpretable model locally around a specific prediction, explaining how the input data contributed to the outcome. 3. **Partial Dependence Plots**: … [189 words] & partial & 1/4 & IBM's AI Explainability toolkit is one of the most mature offerings on the market. It offers global and local explainability, fairness indicators, and counterfactual explanations. [25 words] & correct & 4/4 \\
qwen3-8b & qwen3-8b\_\_ep3 & test\_trained\_exact & ai-in-higher-education-a-bibliometric-analysis:0019:q4 & factual & equity\_accessibility & ai-in-higher-education-a-bibliometric-analysis & 1.000000 & True & What is one cause of bias in AI systems that some researchers warn about? & Some researchers warn that biases in AI can arise from a lack of data diversity. & One cause of bias in AI systems that some researchers warn about is **biased training data**, which can reflect historical prejudices or inequalities present in the data sources. [28 words] & partial & 0/1 & Some researchers warn that biases in AI can arise from a lack of data diversity. [15 words] & correct & 1/1 \\
gemma-4-e4b-it & gemma-4-e4b-it\_\_ep3 & test\_trained\_exact & innovative-educational:0063:q2 & definition & assessment & innovative-educational & 0.600000 & True & Which core design principles does the effectiveness of personalized assessment in higher education depend on? & The effectiveness of personalized assessments hinges critically on the thoughtful integration of core design principles. These include timely and actionable formative feedback, the sophisticated use of adaptive technologies and intelligent tutoring systems, and the ethical application of learning analytics to inform tailored pathways. & The effectiveness of personalized assessment in higher education fundamentally depends on several core design principles: 1. **Alignment with Learning Outcomes (Constructive Alignment):** Assessments must directly measure the specific knowledge, skills, and competencies defined in the learning objectives. If the assessment doesn't align with what students are supposed to learn, personalization is irrelevant. 2. **Diagnostic Capability:** The assessment must be designed not just to grade, but to *diagnose* the student's current … [218 words] & partial & 2/5 & The effectiveness of personalized assessment is contingent upon the thoughtful integration of core design principles. These include timely and actionable formative feedback, the sophisticated use of adaptive technologies and intelligent tutoring systems, and the ethical application of learning analytics to inform tailored pathways. [43 words] & correct & 5/5 \\
qwen3-8b & qwen3-8b\_\_ep3 & test\_trained\_exact & ai-frameworks:0019:q4 & factual & governance & ai-frameworks & 1.000000 & True & What is one instance of an emergent capability revealed when evaluating a newly trained AI model that could lead to harm if left unchecked? & A model capable of autonomously solving CAPTCHA challenges can be used by bad actors to overcome spam protections. Evaluating a newly trained model can reveal such emergent capabilities, which could lead to harm if left unchecked. & One instance of an emergent capability that could lead to harm is **unintended bias amplification**. When evaluating a newly trained AI model, it may inadvertently reinforce or amplify existing biases present in the training data, leading to unfair or discriminatory outcomes in applications such as hiring, lending, or law enforcement. If left unchecked, this can perpetuate systemic inequities and harm marginalized groups. [62 words] & incorrect & 0/1 & A model capable of autonomously solving CAPTCHA challenges can be used by bad actors to overcome spam protections. Evaluating a newly trained model can reveal such emergent capabilities, which could lead to harm if left unchecked. [36 words] & correct & 1/1 \\
\end{longtable}

